\cleardoublepage

\section{ICESat-2卫星与验证数据}

\subsection{ICESat-2 任务概况}
在全球气候变暖背景下，极地冰盖、海冰和积雪覆盖正在发生显著变化。冰盖表面高程变化关系到冰量收支和海平面变化，海冰干舷和海冰厚度则直接影响极区海气相互作用与能量交换。由于极区现场观测受环境、航线和季节限制明显，能够长期、稳定获取全球地表高程信息的星载激光测高任务具有重要价值。ICESat-2 是美国国家航空航天局面向全球地表高程变化监测建设的第二代 Ice, Cloud, and land Elevation Satellite 任务，为冰冻圈变化研究提供了高分辨率沿轨观测基础\citep{ICESat2_MissionOverview}。

ICESat-2 与早期 ICESat 任务之间具有明确的继承关系。2003 年 NASA（National Aeronautics and Space Administration，美国国家航空航天局）发射 ICESat，卫星搭载 GLAS（Geoscience Laser Altimeter System，地球科学激光测高系统），通过星载激光测高方式获取全球尤其是极地冰盖区域的高程观测。ICESat 在轨运行期间形成的长时间序列资料，为冰盖高程变化、海冰干舷和陆地高程研究奠定了重要基础。ICESat-2 于 2018 年 9 月发射，作为 ICESat 的后继星，其任务目标不是简单延续已有观测，而是在采样密度、观测方式和空间覆盖能力上进一步提升全球激光测高能力\citep{Zhu2020_ICESat2_ATLAS_Review,ICESat2_MissionOverview}。

从任务定位看，ICESat-2 面向多类地表目标提供高精度高程观测，主要科学应用包括冰盖高程变化监测、海冰干舷与海冰厚度估算、森林植被高度测量、陆地地形和内陆水体高程观测等。对于本文关注的冰冻圈问题，ICESat-2 的价值不仅在于能够获取雪面或冰面高程，还在于其观测可用于分析海冰表面形态、干舷变化以及积雪对海冰厚度反演的影响。已有研究常将 ICESat-2 与雷达测高、被动微波或机载观测结合，用于改善海冰厚度和积雪参数估计\citep{Kwok2017}。

ICESat-2 的轨道设计也服务于高纬度持续观测需求。卫星运行高度约为 \(500\,\mathrm{km}\)，轨道倾角约为 \(92^\circ\)，观测范围可覆盖约 \(88^\circ\mathrm{S}\sim88^\circ\mathrm{N}\)，重复周期约为 91 天。这样的近极地轨道使其能够反复经过南北极及高纬陆地区域，在多周期观测中积累沿轨高程剖面。对于北极海冰和积雪研究而言，广泛的纬度覆盖有利于获取海冰区不同轨道、不同季节的表面观测；91 天重复周期则为比较同一区域附近的时序变化提供了基础。

与传统全波形激光雷达不同，ICESat-2 搭载的 ATLAS 采用光子计数激光雷达体制。该体制并不直接输出连续回波波形，而是记录返回探测器的离散光子事件，并根据光子往返飞行时间计算测距结果，再结合卫星轨道、姿态和指向信息确定光子的经纬度与高程。ATLAS 工作波长为 \(532\,\mathrm{nm}\)，具有较高脉冲重复频率、较小地面光斑和六束同步观测能力，因此能够在沿轨方向形成密集的光子点云。上述特征使 ICESat-2 在刻画地表细尺度起伏和弱返回信号方面具有不同于传统测高任务的观测优势。

需要指出的是，ICESat-2 官方高级产品主要围绕地表高程、海冰干舷和植被高度等应用生成，并不直接给出积雪深度或雪深直接反演结果。本文关注的不是常规高程产品中的单一表面高程，而是 ATL03 光子点云在雪面附近及雪面以下保留下来的次表层弱信号。激光光子进入雪层后可能经历多次散射，部分返回光子相对主表面回波产生延迟，从而在垂向剖面中表现为次表层长尾信号。若能够将这类弱信号与仪器噪声、背景光和系统响应拖尾区分开，就可能在不依赖无雪期参考高程的情况下，为雪深直接反演提供新的观测信息\citep{Hu2022,Lu2022Unc}。

在 ICESat-2 公开数据产品体系中，ATL03 是全球地理定位光子数据产品，也是许多高级产品生成的基础。ATL03 保留逐光子事件的时间、经纬度和高程等基础信息，使研究者能够根据具体问题重新组织光子点云，而不是只能使用已经参数化后的高程或干舷结果。本文选择 ATL03 作为主反演数据，主要原因是其较完整地保留了雪面附近光子事件的空间和垂向分布，便于构建相对于参考表面的光子剖面，并进一步提取雪层多次散射造成的次表层长尾结构\citep{NSIDC_ATL03}。

\subsection{ATLAS光子计数激光雷达系统}
ATLAS 是 ICESat-2 搭载的 Advanced Topographic Laser Altimeter System，工作波长为 532\,nm。该系统采用光子计数探测方式，即记录单个返回光子的到达时间，并结合卫星轨道、姿态和大气改正信息计算光子地理位置与高程。与传统全波形激光雷达输出连续回波波形不同，ATLAS 输出的是大量离散光子事件；这些事件在沿轨方向和垂向上共同组成光子点云。

ICESat-2 共发射六束激光，组成三组光束对。每组光束对包含一条强束和一条弱束，两者在地面上形成横向相邻的观测轨迹。通常所说的轨道1、轨道2和轨道3表示三组地面轨迹，而不是三次独立卫星过境。强束接收的有效光子数通常更多，表面回波更稳定，也更有利于提取表面以下弱长尾信号；弱束虽然同样具有空间覆盖价值，但在噪声背景较强或雪层信号较弱时更容易出现剖面不稳定问题。

基于上述差异，本文主要使用强束结果进行雪深反演和 OIB 验证统计。第4章定量表格统一采用轨道2强束配对样本，以保证样本定义、日期覆盖和统计口径一致。轨道1和轨道3在部分图件中保留，主要用于观察同一日期内不同轨迹分支的沿轨形态是否一致，而不与轨道2混合计算统一精度指标。

光子计数体制的优势在于灵敏度高，能够保留表面附近和表面以下的弱返回信号。对于普通地表测高任务，研究重点往往是从光子点云中提取主表面高程；而对于本文关注的雪深问题，表面以下的弱长尾并非完全是噪声，其中可能包含由雪层内部多次散射造成的延迟回波。正是由于 ATL03 保留了原始光子事件，本文才能进一步从雪面以下弱光子分布中提取多次散射信息。

需要注意的是，ATLAS 的高灵敏度也意味着观测会受到太阳背景、探测器暗计数、后脉冲和系统响应拖尾等因素影响\citep{Lu2022Unc}。这些因素会改变表面附近和表面以下回波剖面的形态，因此后续方法章节必须对背景噪声和系统响应进行处理，不能直接把原始光子计数剖面等同于雪层真实散射响应。ATLAS 系统结构、光束布局、光子计数流程和主要误差来源如\autoref{fig:atlas-atl03-intro}所示。

\begin{figure}[!htbp]
    \centering
    \includegraphics[width=0.92\textwidth, trim=0 6 4 8, clip]{final/ATL03_intro}
    \caption{\label{fig:atlas-atl03-intro}ATLAS 光子计数激光雷达系统、光束结构与 ATL03 光子事件示意图。（a）卫星观测与测距几何；（b）三组强弱光束对的空间布局；（c）光子计数探测与 ATL03 事件输出流程；（d）光子计数观测优势及主要误差来源。}
\end{figure}
\FloatBarrier

\subsection{ATL03 光子数据产品介绍}
ATL03 是 ICESat-2 卫星搭载的 ATLAS 光子计数激光雷达生成的 L2A（Level-2A，二级 A 类数据产品）全球地理定位光子数据产品，也是 ICESat-2 公开科学数据中最接近原始观测、最基础的逐光子点云产品\citep{NSIDC_ATL03}。相较于已经提取地表高程、海冰干舷或植被高度等物理量的高级产品，ATL03 更接近 ATLAS 的原始观测形态：它以单个光子事件为基本记录对象，在遥测数据、精密轨道、姿态信息和光斑定位模型的基础上，经过必要的大气、潮汐等地球物理改正后，给出每个光子的地理位置和高程信息。因此，ATL03 是后续多类 ICESat-2 高级产品生成的基础，也是研究者重新组织光子点云并开展自定义算法处理的主要入口。

ATL03 数据以 HDF5（Hierarchical Data Format version 5，第五版层次化数据格式）文件形式发布，文件名中通常包含观测日期与时刻、参考地面轨道编号、周期号、轨道段号以及版本修订号等信息。这样的组织方式使每个数据文件都可以与具体过境时间、轨道位置和重复周期对应起来，便于按日期、轨道段和研究区域筛选数据。对于本文的北极海冰和陆地区域实验而言，文件层面的时间和轨道信息主要用于确定候选过境轨段，而具体的雪深反演仍需要进入逐光子层面重新构建垂向剖面。

从变量内容看，ATL03 为每个光子事件保存了经度、纬度、高程、接收时间等基本地理定位信息，同时还提供沿轨距离、光束编号、背景计数和信号置信度等辅助信息。这些字段并不是彼此独立的参数清单，而是共同描述光子在时间、空间和垂向上的位置关系。经纬度用于研究区裁剪和外部雪深数据配对，高程用于识别主表面并定义相对深度，接收时间和沿轨距离用于将离散光子排列成沿轨序列，光束编号则用于区分不同强弱束观测。背景计数和信号置信度虽然不能直接给出雪深，但可以为质量控制和噪声解释提供依据。

对于常规高程应用，ATL03 往往只是生成 ATL06（ICESat-2/ATLAS Level-3A Land Ice Height，ICESat-2/ATLAS 三级 A 类陆冰高程产品）、ATL07（ICESat-2/ATLAS Level-3A Sea Ice Height，ICESat-2/ATLAS 三级 A 类海冰高程产品）等高级产品的中间基础数据，研究者可以直接使用后者提供的陆冰高程或海冰参数。然而，本文关注的是雪面附近和雪面以下的次表层弱信号，这类信号在高级产品的参数化处理过程中可能被弱化或不再完整保留。ATL03 保留逐光子事件的离散分布，使本文能够围绕主表面重新对齐光子高程，并在统一参考表面下统计雪面以下的次表层长尾信号。

因此，本文选择 ATL03 作为雪深直接反演的主输入数据，而不是直接使用已经生成的高程或干舷产品。具体处理时，本文首先利用 ATL03 光子事件建立沿轨时间--高程分布，识别局地主表面并构建相对深度坐标；随后在沿轨窗口内聚合光子点云，形成可用于背景扣除、系统响应去卷积和多次散射统计矩计算的光子剖面。由此可见，ATL03 在本文中不仅是数据来源，更是连接 ICESat-2 光子计数观测与雪深反演模型的基础数据载体。

\subsection{本文使用的数据与验证方案}
\subsubsection{数据体系与研究区域}
本文的数据体系围绕``反演输入--海冰验证--陆地对照''三类功能组织。反演输入为 ICESat-2 ATL03 光子数据及其沿轨雪深反演结果，核心作用是提供光子计数观测和雪面以下弱长尾信号；海冰验证数据为 Operation IceBridge（OIB）机载雪深观测，用于检验 ICESat-2 反演结果在北极海冰局地重叠航段上的一致性；陆地对照数据包括 University of Arizona（UA）4\,km 日尺度雪深产品和 Canadian Meteorological Centre（CMC）24\,km 逐日雪深分析产品，用于说明当前流程在陆地格网雪深背景下的响应差异和适用边界。这样的组织方式把数据来源、地表类型和验证意义放在同一框架下，避免将不同性质的参考数据简单合并为统一精度评价。

从研究区域看，OIB 验证样本位于北极海冰区，选取 2019 年 4 月 12 日、4 月 19 日和 4 月 22 日三次与 ICESat-2 轨段形成近时空交会的机载航线；UA 和 CMC 对照样本位于 2020 年 3 月 27 日 ICESat-2 陆地轨段附近，主要覆盖北美陆地区域。两类区域的地表条件和验证含义明显不同：海冰样本以相对平坦的冰雪界面和高密度机载航线为主，更适合开展沿轨局地配对；陆地样本受到地形起伏、植被、格网平均和产品融合方式影响，更适合作为方法迁移边界的说明。因此，本文在第4章中分别统计 OIB 海冰验证结果和 UA/CMC 陆地对比结果，不将二者合并为一个总体指标。

主要数据及用途如\autoref{tab:data-source-summary}所示。该表并不只是列出数据名称，而是说明每类数据在本文流程中的位置：ATL03 对应反演模型输入，OIB 对应独立机载验证，UA 和 CMC 对应陆地格网参考。后续各小节按照这一顺序分别说明数据内容、变量特征、使用方式和验证限制。

\begin{table}[htbp]
    \centering
    \caption{\label{tab:data-source-summary}本文使用的主要数据}
    \small
    \begin{tabularx}{0.98\linewidth}{p{0.20\linewidth}p{0.20\linewidth}p{0.28\linewidth}X}
        \toprule
        数据 & 作用 & 时间 & 空间范围 \\
        \midrule
        ICESat-2 ATL03 & 雪深反演数据 & 2019 年 4 月 12、19、22 日；2020 年 3 月 27 日 & 北极海冰、美国本土陆地区域 \\
        OIB 雪深数据 & 海冰验证数据 & 2019 年 4 月 12、19、22 日 & 北极海冰 \\
        UA 4 km 雪深 & 陆地验证数据 & 2020 年 3 月 27 日 & 美国本土陆地区域 \\
        CMC 雪深分析 & 陆地格网验证数据 & 2020 年 3 月 27 日 & 北半球陆地区域 \\
        \bottomrule
    \end{tabularx}
\end{table}

数据来源和研究区域之间的关系如\autoref{fig:data-region-overview}所示。该图将 ATL03 反演数据、OIB 海冰验证样本以及 UA/CMC 陆地验证样本放在同一框架下，用于说明不同数据在空间范围、时间和地表类型上的分工。需要注意的是，该图表达的是数据体系和研究区域的对应关系，而不是全部原始观测点的完整空间分布；具体轨迹交会、配对子集和结果统计将在第4章展开。

\begin{figure}[!htbp]
    \centering
    \includegraphics[width=0.8\textwidth, trim=10 10 8 12, clip]{final/data_compare}
    \caption{\label{fig:data-region-overview}本文使用的数据来源、研究区域与验证场景示意图。（a）ATL03、OIB、UA 和 CMC 数据的功能分工；（b）北极海冰验证区域与北美陆地对照区域的空间分布。蓝色轨迹表示 ICESat-2，橙色航线表示 OIB，网格背景表示 UA/CMC 雪深产品。}
\end{figure}
\FloatBarrier

\subsubsection{ICESat-2 ATL03反演数据}
ICESat-2 ATL03 是本文雪深直接反演的核心输入数据，也是连接星载光子计数观测和雪深反演模型的基础数据载体\citep{NSIDC_ATL03}。与 ATL06、ATL07 等已经提取高程或海冰参数的高级产品相比，ATL03 保留逐光子事件的经纬度、高程、接收时间、沿轨位置、光束编号、背景计数和信号置信度等信息，允许研究者根据具体问题重新构建光子剖面。本文关注的不是单一雪面高程，而是雪面附近及雪面以下由多次散射可能形成的弱长尾信号，因此必须使用 ATL03 这类逐光子产品，而不能直接从高级产品中读取一个已经参数化的雪深量。

在本文实验中，ATL03 数据主要对应两类场景。第一类为 2019 年 4 月北极海冰样本，分别与 OIB 的 4 月 12 日、4 月 19 日和 4 月 22 日航线形成局地交会，用于海冰雪深反演验证；第二类为 2020 年 3 月 27 日陆地轨段，用于与 UA 和 CMC 格网雪深数据开展对照。两类场景都使用相同的基本 ATL03 变量读取方式，但解释重点不同：海冰样本强调与机载 OIB 航线的局地近时空配对，陆地样本强调同一沿轨反演流程离开海冰环境后的响应特征。

具体处理时，ATL03 光子事件首先按日期、轨道段和光束分支组织。每个候选轨段保留光子时间、高程、经纬度、沿轨距离、光束编号、背景计数和质量标识，并依据研究区域进行空间裁剪。由于强束通常具有更高的有效光子数量和更稳定的表面回波，本文在定量验证中优先使用强束结果；轨道1、轨道2和轨道3可用于展示不同分支的沿轨形态，但第4章 OIB 定量表格统一采用轨道2强束配对样本，以保证样本定义和统计口径一致。

ATL03 数据在本文中还承担中间处理链条的组织作用。反演流程会将逐光子事件转换为时间--高程剖面，识别局地主表面后再转为相对深度坐标，并在滑动沿轨窗口内聚合得到后向散射剖面。随后，剖面进入背景噪声扣除、系统响应去卷积和路径长度统计矩计算。也就是说，ATL03 并不是一个只在数据表中出现的外部输入，而是后续所有物理反演步骤的起点；对其变量选择、光束筛选和轨段组织方式的设定，会直接影响剖面信噪比、主表面识别稳定性以及最终雪深结果。

\subsubsection{OIB海冰雪深验证数据}
OIB 海冰雪深数据是本文最主要的独立验证资料。Operation IceBridge 是 NASA 在 ICESat 与 ICESat-2 任务之间及极地观测需求下开展的机载观测计划，其海冰产品包含沿飞行航线获取的自由板、雪深、海冰厚度、表面粗糙度及不确定度等信息\citep{NSIDC_IDCSI4}。本文使用的 OIB 数据来自 IceBridge L4（Level-4，四级数据产品）Sea Ice Freeboard, Snow Depth, and Thickness 产品，并围绕 2019 年 4 月 12 日、4 月 19 日和 4 月 22 日三次北极海冰航线组织。这三天均存在可与 ICESat-2 ATL03 反演结果进行近时空比较的局部重叠航段，因此适合用于检验本文方法在海冰场景下恢复背景雪深量级的能力。

从文件组织看，OIB 数据既保留原始文本形式，也整理为更适合验证的轻量化雪深表。原始文件中包含纬度、经度、海冰厚度、厚度不确定度、平均自由板、ATM（Airborne Topographic Mapper，机载地形测绘仪）自由板、自由板不确定度、雪深、雪深不确定度和表面粗糙度等字段，缺测值通常以 \(-99999\) 表示，经度也可能采用 \(0\sim360^\circ\) 表达。为了与 ICESat-2 结果匹配，本文优先使用整理后的 \texttt{oib.csv} 文件，其中保留 \texttt{lat}、\texttt{lon}、\texttt{snow\_depth} 和 \texttt{snow\_depth\_unc} 等直接参与验证的变量，并将经度转换到 \([-180^\circ,180^\circ]\) 范围。三期 OIB 简化表的有效雪深样本数分别约为 32579、26778 和 16921 条，能够提供高密度沿航线参考。

OIB 数据的优势在于采样密度高、观测对象明确，能够在局地交会区域内给出比格网产品更接近沿轨尺度的雪深参考。对于本文方法而言，这一点尤为重要：ICESat-2 雪深反演结果本身沿轨分布，且每个反演点都由局地光子剖面统计得到；如果使用过粗的参考数据，反演误差会与尺度差异混在一起。OIB 沿航线数据虽然仍然不是与 ICESat-2 完全同足迹的观测，但其空间采样方式更接近沿轨验证需求，因此本文将其作为海冰雪深反演精度评价的主要依据。

同时，OIB 数据的使用边界也需要明确。三期航线只覆盖北极海冰区的局部飞行轨迹，且与 ICESat-2 地面轨迹之间是局部交会关系，并非整条航线完全重合。最终验证样本只代表能够满足空间邻域和沿轨窗口约束的局部重叠航段，而不能直接代表同日整个 OIB 航线或整个北极海冰区的雪深场。后文的 Bias、RMSE 和相关系数等指标都应在这一配对子集口径下解释。

\subsubsection{UA原位同化陆地雪深数据}
UA 雪深数据属于原位约束较强的陆地雪深参考资料，本文使用其 daily 4 km gridded snow depth 产品开展陆地对照\citep{UA_SnowDepth}。该产品以日尺度格网形式提供雪深和雪水当量等变量，数据构建过程综合了 SNOTEL（Snow Telemetry，雪遥测网络）、COOP（Cooperative Observer Program，合作观测员计划）等地面原位观测以及 PRISM（Parameter-elevation Regressions on Independent Slopes Model，独立坡向参数-高程回归模型）温度、降水等格网气象资料。相较于单纯的再分析或业务化大尺度雪深场，UA 产品在美国本土区域受到更多地面站点观测约束，因此更适合用于检查 ICESat-2 沿轨反演结果是否能够反映陆地积雪背景的空间变化。

本文选取 2020 年 3 月 27 日 ICESat-2 陆地轨段与 UA 数据进行对比。该日期对应的 UA 产品提供 4\,km 日尺度格网雪深值，空间分辨率高于 CMC 数据，能够在一定程度上保留局地积雪差异。将 ICESat-2 反演结果投影到 UA 格网后，可以得到每个沿轨反演点附近的格网雪深参考，从而比较星载光子剖面反演结果与原位同化型陆地雪深背景之间的量级和沿轨变化关系。

UA 数据的验证意义与 OIB 不同。OIB 是沿航线机载观测，能够与 ICESat-2 轨迹在局部范围内形成较直接的沿轨比较；UA 则是格网产品，其 4\,km 空间分辨率远大于 ATLAS 光斑和本文局地剖面窗口。一个 UA 格网值代表的是格网范围内经过观测同化和空间插值后的日尺度平均状态，而不是 ICESat-2 足迹处的瞬时雪深。因此，ICESat-2 与 UA 的差异同时包含反演误差、格网平均误差、地形和植被影响、观测时间差以及参考产品本身的不确定性。

基于上述特点，本文不把 UA 对比解释为陆地雪深最终精度检验，而是将其作为方法适用性边界分析的一部分。如果 ICESat-2 反演雪深与 UA 在沿轨变化上具有一致趋势，说明多次散射剖面中可能保留了与陆地积雪背景相关的信息；如果二者误差明显大于 OIB 海冰样本，则说明围绕海冰场景设定的主表面识别、窗口尺度和参数体系不能直接迁移到复杂陆地场景。这样的解释比简单给出一个统一 RMSE 更符合两类数据的观测尺度差异。

\subsubsection{CMC陆地格网雪深分析数据}
CMC 雪深数据为 Canadian Meteorological Centre 逐日雪深分析产品，本文将其作为陆地格网雪深对照资料\citep{CMC_SnowDepth}。该产品覆盖北半球陆地区域，空间分辨率约为 24\,km，提供逐日雪深分析结果。与 UA 产品相比，CMC 的空间分辨率更粗，场的平滑程度更高，更适合反映大尺度陆地积雪背景，而不适合解释短距离内的局地雪深峰谷。

本文使用 2020 年 3 月 27 日的 CMC 雪深分析结果，与同日 ICESat-2 陆地轨段进行配准比较。由于 CMC 覆盖范围广、格网连续性较强，沿轨 ICESat-2 结果通常能够找到对应格网值，这使其可以作为 UA 对比之外的另一种陆地参考。UA 和 CMC 同时使用的目的，是区分两类陆地参考资料的尺度和构建方式差异：UA 更接近站点约束下的 4\,km 局地背景，CMC 则代表较粗分辨率、业务化的大尺度积雪分析场。

CMC 数据的优势在于空间覆盖稳定、时间连续且数据结构规则，便于与 ICESat-2 沿轨结果进行批量匹配。但这种优势也伴随明显的尺度限制。24\,km 格网会平滑山地、林地和局地风成积雪等细尺度变化，一个格网值可能包含多种地表类型和积雪状态。当 ICESat-2 轨段穿过地形起伏明显或积雪空间异质性较强的区域时，沿轨反演点与 CMC 格网值之间的差异不能简单归结为反演算法误差，也可能来自格网代表性不足。

因此，本文将 CMC 数据定位为大尺度陆地背景对照，而不是高精度逐点验证基准。若 ICESat-2 反演曲线与 CMC 在区域雪深量级上相近，可以说明当前结果与大尺度积雪背景没有明显矛盾；若二者在局地细节上差异较大，也不能直接判定反演失败。将 CMC 与 UA 并列使用，有助于在第4章中说明参考数据尺度越粗、地表异质性越强时，ICESat-2 雪深反演结果的解释就越依赖配对口径和场景条件。

\subsubsection{验证配对与预处理口径}
本文的验证配对以地表类型分开进行。海冰部分以 ICESat-2 反演雪深点为基本单位，对每一个 ICESat-2 点，在 OIB 航线中选取满足 4\,km 空间邻域约束且落入 350\,m 沿轨窗口的有效 OIB 雪深观测，并以窗口内 OIB 雪深平均值作为该点的机载参考值。若窗口内没有有效 OIB 观测，则该 ICESat-2 点不进入统计。因此，第4章中的 OIB 配对数表示``ICESat-2 反演点--OIB 邻域平均值''形成的有效样本对数量，而不是原始 OIB 采样点数量。跨日期统计统一使用轨道2强束结果，以避免不同轨道分支混合导致样本重复或统计口径不一致。

陆地部分则将 2020 年 3 月 27 日 ICESat-2 沿轨反演点分别与 UA 和 CMC 当日格网雪深进行匹配。由于 UA 和 CMC 均为日尺度格网产品，本文不再进行小时级时间插值，而是以同日格网值作为参考背景；空间上根据反演点经纬度提取对应或最近格网单元的雪深值，并记录无效格网或缺测情况。该配对方式适合评价沿轨结果与陆地雪深背景的相对一致性，但不具备 OIB 那样的高密度沿航线独立验证含义。

在进入反演和验证之前，各类数据均进行基础预处理。ATL03 数据读取光子时间、高程、经纬度、沿轨距离、光束编号、背景计数和质量信息，并按研究区域、轨道段、强弱束和有效光子数量进行筛选；OIB 数据剔除缺测雪深，统一经度表达，并保留雪深及其不确定度字段；UA 和 CMC 数据统一日期、格网坐标和雪深变量，转换为可与 ICESat-2 沿轨结果配对的形式。本节只说明数据来源、观测变量和配对口径，主表面识别、背景扣除、系统响应去卷积、路径长度统计矩和雪深迭代求解属于方法内容，将在后续方法小节集中展开。

\cleardoublepage

\section{基于多次散射的雪深反演方法}

本章建立从 ICESat-2 ATL03 光子点云到沿轨雪深结果的完整反演方法。整体流程如\autoref{fig:method-workflow}所示：首先对 ATL03 光子事件进行筛选，并通过主表面识别与表面对齐建立统一深度坐标；随后将离散光子投影为时间--深度剖面并沿轨滑动聚合；在此基础上扣除背景噪声、校正系统响应拖尾，得到用于物理反演的校正后向散射剖面；最后将剖面长尾转化为路径长度分布及其统计矩，并通过吸收、散射参数迭代求解雪深。

\begin{figure}[!htbp]
    \centering
    \includegraphics[width=0.94\textwidth, trim=8 8 12 12, clip]{final/method/snow_retrieval_pipeline}
    \caption{\label{fig:method-workflow}基于 ATL03 光子剖面和多次散射统计矩的雪深反演流程图。（1）ATL03 输入与预处理；（2）剖面处理与系统响应校正；（3）多次散射路径长度统计与雪深求解；（4）OIB、UA 和 CMC 结果评价。}
\end{figure}
\FloatBarrier

\subsection{雪层多次散射理论与雪深反演原理}
在 532\,nm 波段下，干雪通常可近似视为弱吸收、强散射的颗粒介质。雪层并不是一个光学均匀的连续体，而是由大量冰晶颗粒、孔隙空气和局地层理共同组成的多相介质。入射激光到达雪面后，一部分光子在空气--雪界面附近被快速反射，形成 ATL03 剖面中的主表面回波；另一部分光子进入雪层内部，在雪粒与孔隙之间传播。由于冰粒折射率与空气不同，光子每遇到一个雪粒或粒间界面，都可能发生折射、反射或散射，传播方向随之改变。相对于吸收过程，干雪在可见光波段的散射作用更强，因此进入雪层的光子往往不会只沿垂直方向完成一次简单往返，而是在多个雪粒之间经历一系列随机方向改变后才重新离开雪面，或在传播过程中被冰晶吸收而消失。

从单个光子的角度看，雪层内传播可理解为由许多短步长组成的随机游走过程。光子进入雪层后，先在近表层发生一次或少数几次散射；若仍未逃逸出雪层，则会继续向下或斜向传播，并在后续雪粒处再次散射。每一次散射都会改变光子的传播方向和后续路径，使其实际行进距离明显大于从卫星到雪面再返回的几何路径。部分光子在经历若干次散射后重新朝上并离开雪层，被 ATLAS 探测器记录为返回光子；另一部分光子则可能继续向更深处传播、被吸收，或因多次散射后不再回到接收视场内。由于每个光子的散射次数、步长、方向序列和吸收概率均不同，返回光子不再对应单一固定路径，而是形成具有统计分布的路径长度集合。

这种路径延长会直接改变光子计数激光雷达观测到的垂向剖面形态。若地表为理想硬目标，返回光子主要集中在系统脉冲宽度控制的窄时间范围内，表面以下不会出现显著的物理长尾。对于积雪表面，进入雪层并多次散射后返回的光子具有更长的传播时间，因而在以主表面为零点的深度坐标中表现为表面以下的延迟回波。散射次数较少、路径较短的光子集中在浅层附近；散射次数较多、路径较长的光子对应更深的等效路径位置，并形成逐渐衰减的次表层尾部。雪层越厚，光子在逃逸前可探索的散射空间越大，长路径光子的出现概率通常越高，剖面尾部也更容易展宽；雪粒尺度、密度、层理和吸收系数则通过改变平均自由程、散射方向分布和长路径光子的存活概率，进一步调节尾部衰减速度。因此，多次散射并不只是增加回波噪声，而是在剖面形态中留下了与雪层光学厚度和物理厚度相关的统计信息。

需要区分的是，ATL03 剖面中的表面以下坐标并不等同于光子实际到达的几何深度。经过表面对齐后，深度坐标 \(z\) 表示相对于主表面的双程传播延迟；对于后向返回光子，可用等效路径长度
\begin{equation}
    \langle L\rangle=2z
\end{equation}
将垂向延迟与往返光程联系起来。其中，\(\langle L\rangle\) 表示返回光子的平均等效路径长度，\(L\) 表示单个光子被观测延迟投影后的等效路径长度，\(z\) 表示相对于主表面的表面以下深度。真实光子路径在雪层内是折线式的三维随机轨迹，正是这种等效路径长度的统计分布，使光子计数剖面能够从时间延迟观测转化为雪层内部多次散射过程的统计描述。光子在雪层中多次改变方向并形成等效长路径的概念如\autoref{fig:pathlength-distribution}所示。

\begin{figure}[!htbp]
    \centering
    \includegraphics[width=0.58\textwidth]{final/method/pathlength_distribution_scheme}
    \caption{\label{fig:pathlength-distribution}均匀照明条件下光子扩散平均行进路径示意图\citep{blanco2003invariance}}
\end{figure}
\FloatBarrier

由图~\ref{fig:pathlength-distribution}可见，光子在散射介质中的实际传播轨迹并不是连接入射点与出射点的直线，而是在介质内部经历多次方向改变后形成曲折路径；当平均自由程远小于介质尺度时，单个光子的实际行进距离可以显著大于几何穿透深度。因此，ATL03 剖面中的表面以下延迟信号更适合解释为路径长度统计结果，而不是光子到达某一固定几何深度后的直接反射。

扩散理论与蒙特卡洛辐射传输模拟表明，在一定的雪层均匀性和弱吸收假设下，雪层后向返回光子的路径长度分布可近似用 Gamma 型分布描述，其低阶统计矩与雪深之间存在稳定关系\citep{Hu2022}。这一结论的物理含义是：雪深控制了光子在雪层中可经历的随机游走空间尺度，而路径长度分布的平均值和离散程度则记录了光子在该空间中传播和逃逸的统计结果。对于较浅雪层，光子较容易在少量散射后返回，路径长度分布集中于短路径；对于较深雪层，光子可能在更厚介质内经历更长扩散过程，长路径尾部占比增加，平均路径长度随之增大。因此，路径长度统计特征与雪深之间并非纯经验相关，而是由雪层作为强散射介质时的扩散传播机制决定。

在实际反演中，首先需要从观测剖面中提取代表雪内多次散射的校正后向散射信号。后文通过背景扣除和系统响应去卷积得到的剖面记为 \(\beta(z)\)。由于光子在雪内传播时会受到吸收衰减，直接使用 \(\beta(z)\) 会低估长路径光子的相对贡献，因此需对吸收效应进行校正。设路径长度概率密度为 \(p(L)\)，雪层吸收系数为 \(k_a\)，则可将去卷积后的后向散射剖面归一化为
\begin{equation}
    p(L)=\frac{\beta(L/2)e^{k_aL}}{\int_{0}^{\infty}\beta(\xi/2)e^{k_a\xi}\,\mathrm{d}\xi}.
\end{equation}
将变量换为深度坐标 \(z\) 后，雪层有效反照率可表示为
\begin{equation}
    a=\frac{\int_{0}^{\infty}\beta(z)\,\mathrm{d}z}{\int_{0}^{\infty}\beta(z)e^{2k_az}\,\mathrm{d}z},
\end{equation}
路径长度的 \(n\) 阶矩为
\begin{equation}
    \langle L^n\rangle =
    \frac{\int_{0}^{\infty}(2z)^n\beta(z)e^{2k_az}\,\mathrm{d}z}
    {\int_{0}^{\infty}\beta(z)e^{2k_az}\,\mathrm{d}z}.
\end{equation}
其中，\(p(L)\) 表示等效路径长度为 \(L\) 时的概率密度，\(\beta(\cdot)\) 为校正后向散射剖面在对应深度处的强度，\(k_a\) 为雪层吸收系数，\(\xi\) 为路径长度积分变量，\(a\) 为由吸收校正前后剖面积分比得到的有效反照率，\(\mathrm{d}\xi\) 和 \(\mathrm{d}z\) 分别为路径长度和深度积分微元。\(z\) 为相对主表面的深度坐标，\(\langle L^n\rangle\) 为路径长度的 \(n\) 阶矩，\(n\) 表示矩的阶数；吸收校正因子 \(e^{2k_az}\) 用于补偿长路径光子更容易被吸收的事实，使矩计算更接近无吸收条件下的路径长度统计。

根据多次散射理论，路径长度一阶矩与雪深 \(H\) 之间满足
\begin{equation}
    H=\frac{\langle L\rangle}{2}.
\end{equation}
其中 \(H\) 为待反演的雪深，\(\langle L\rangle\) 为路径长度一阶矩。该式说明，平均等效路径长度可作为雪层厚度的直接统计量：当雪层变厚时，返回光子的平均传播延迟随之增加，反演得到的 \(H\) 也相应增大。二阶矩和三阶矩在理论上也可提供雪深约束，即
\begin{equation}
    H=\left(\frac{\langle L^2\rangle}{k_{sd}}\right)^{1/3},
    \qquad
    H=\left(\frac{\langle L^3\rangle}{k_{sd}^2}\right)^{1/5},
\end{equation}
其中 \(k_{sd}\) 为雪层扩散散射系数，\(\langle L^2\rangle\) 和 \(\langle L^3\rangle\) 分别为路径长度的二阶矩和三阶矩，\(H\) 仍表示雪深。二阶矩描述路径长度的离散程度，可用于约束扩散散射系数和吸收系数迭代；三阶矩对长路径尾部更加敏感，理论上能够补充深层传播信息，但在实际 ATL03 剖面中也更容易受到深层弱信号、背景扣除误差和去卷积振荡影响。因此，本文将一阶矩作为雪深估计的主控量，使用二阶矩参与参数迭代，而不把三阶矩作为最终输出的主控量。

上述公式给出了雪深反演的理论核心，但实际 ATL03 剖面并不是雪层散射响应的直接观测。表面以下的观测剖面同时受到太阳背景、随机暗计数、探测器后脉冲和系统响应拖尾影响，可概念化表示为
\begin{equation}
    S_o(z)=R(z)\otimes \beta(z)+b(z)+\varepsilon(z),
\end{equation}
其中 \(S_o(z)\) 为观测剖面，\(\beta(z)\) 为期望恢复的雪层后向散射剖面，\(R(z)\) 为系统响应函数，\(b(z)\) 为局地背景项，\(\varepsilon(z)\) 为随机噪声。因而，路径长度统计矩不能直接由原始光子计数计算，而需要先经过主表面识别、背景扣除、系统响应校正和有效剖面筛选。

综上，基于雪层多次散射的雪深反演成立于以下物理链条：激光进入雪层后在雪粒之间发生多次散射和弱吸收，导致部分返回光子的传播路径被显著延长；路径延长表现为 ATL03 表面对齐剖面中的次表层长尾；长尾经背景扣除、系统响应校正和吸收校正后，可转化为等效路径长度分布；该分布的低阶统计矩由雪层厚度和光学参数共同控制，因此能够用于估计雪深。与依赖无雪期高程或外部数字高程模型的方法不同，该理论直接利用雪层内部多次散射信号本身进行反演，适合构建面向 ICESat-2 光子计数观测的雪深直接估计框架。同时，该方法也依赖干雪弱吸收、主表面识别稳定、背景噪声可扣除以及系统响应可校正等前提；若表面粗糙度、大气散射、融雪或复杂层理显著改变剖面长尾，则反演结果需要结合质量控制和误差分析谨慎解释\citep{Lu2022Unc}。

\subsection{ATL03 光子点云预处理与剖面构建}
\label{sec:atl03-profile-build}
ATL03 光子点云处理对应\autoref{fig:method-workflow}中的观测样本整理、主表面识别和剖面构建环节，其输入为逐光子事件集合，输出为相对于局地主表面的沿轨后向散射剖面 \(S_m(z)\)。这一部分的目标不是直接求雪深，而是将离散、含噪的光子点云转换为具有统一深度零点和稳定信噪比的剖面序列，为后续背景扣除、系统响应校正和路径长度统计提供输入。

\subsubsection{光子数据预处理}
ATL03 光子数据包含地表、雪层内部、大气背景和随机噪声等多种来源的光子事件。若直接对全部光子进行剖面统计，少量高空云散射光子、异常高程点或低质量轨段都可能影响主表面定位，并进一步改变表面以下弱长尾的形态。因此，在构建后向散射剖面之前，需要先对 ATL03 点云进行预处理，使进入统计的光子集合在空间范围、光束属性、质量标识和高程分布上具有一致口径。

本文首先根据 ATL03 文件的观测日期、轨道段和光束编号组织候选数据。ICESat-2 每组光束对包含强束和弱束，强束通常具有更高的返回光子数和更稳定的主表面回波，更适合提取表面以下弱信号。因此，本文在定量反演和验证中优先使用强束光子，并结合航天器飞行方向确定实际对应的左右光束。对于每个候选轨段，保留光子接收时间 \(t_i\)、高程 \(h_i\)、纬度 \(\varphi_i\)、经度 \(\lambda_i\)、背景计数和质量标识等变量，为后续沿轨分段、表面识别和噪声估计提供基础。

在空间筛选方面，本文依据研究区经纬度范围对光子进行裁剪，仅保留落入目标海冰或陆地样本区域的光子。空间裁剪的目的在于避免同一轨段中不同地表类型混入同一剖面统计窗口，尤其是海冰样本中需要尽量保持局地表面性质一致。对于跨越区域边界或明显包含复杂地物混合的轨段，只保留满足研究区约束的连续部分；若裁剪后有效光子数量不足，则该轨段不再进入后续剖面构建。

质量控制主要围绕光子级和轨段级两类异常展开。光子级筛选剔除 ATL03 中质量标识异常、缺少必要定位变量或明显偏离局地高程范围的离群光子。轨段级筛选则将光子按沿轨时间分为短时间段，检查每段内的有效光子数量和垂向分布范围；若某一时间段光子过少，或高程起伏明显超出局地主表面可能范围，则认为该段主表面不稳定或受云、大气散射及其他噪声影响较强，不参与后续统计。该处理与参考文献中先剔除陆地、大气噪声再统计信号廓线的思路一致，其核心是保证剖面统计前的光子样本能够代表目标表面附近的有效返回信号。

完成上述筛选后，还需要进行高程粗筛和时间标准化。高程粗筛以候选轨段的局地中位高程或初始表面高度为中心，保留一定高度范围内的光子，以削弱远离表面的高空背景光子和深层异常点对二维直方图的影响。时间标准化则以轨段起始时刻为参考，将绝对接收时间转换为相对时间 \(t_r=t_i-t_0\)，使光子能够按照沿轨顺序进入后续分箱统计。经过预处理后得到的有效光子集合可表示为
\begin{equation}
    \mathcal{P}=\{(t_{r,i},h_i,\varphi_i,\lambda_i)\mid i=1,2,\ldots,N\},
\end{equation}
其中 \(\mathcal{P}\) 为预处理后的有效光子集合，\(t_{r,i}\) 为第 \(i\) 个光子的相对接收时间，\(h_i\) 为光子高程，\(\varphi_i\) 和 \(\lambda_i\) 分别为纬度和经度，\(N\) 为通过筛选的有效光子数。该集合保留了构建沿轨剖面所需的时间、空间和高程信息，同时尽量削弱非目标光子和异常轨段对主表面识别及多次散射信号提取的影响。

\subsubsection{主表面识别与表面对齐}
主表面识别与表面对齐是 ATL03 光子剖面构建中的关键步骤，决定了深度坐标的零点。如果主表面零点识别错误，后续所有深度、路径长度和雪深都会产生系统偏差。本文采用基于时间和高程二维统计直方图的主表面识别方法：先将预处理后的光子时间转换为相对时间 \(t_r=t_i-t_0\)，再在 \((t_r,h)\) 空间中构建二维直方图。对于每一个时间分箱，计数最大的高程位置通常对应最强的表面回波，可作为初始表面估计值。由于原始峰值序列中仍可能存在跳变和局部异常，本文进一步通过缺测插值、异常值剔除和局部回归平滑获得连续表面 \(h_s(t)\)。

如\autoref{fig:surface-tracking-dates}所示，2019 年 4 月 19 日典型轨段的参考表面沿高密度光子分布稳定展开。灰色点云表示原始 ATL03 光子分布，红色轨迹表示识别出的主表面。该图用于说明表面识别是否稳定，而不是直接展示雪深结果。

\begin{figure}[!htbp]
    \centering
    \includegraphics[width=0.92\textwidth]{final/processing/0419_Day_profile_strg3_ATL03_20190419131031_03250304_007_01.surface.jpg}
    \caption{\label{fig:surface-tracking-dates}2019 年 4 月 19 日典型 ATL03 轨段的主表面识别结果。灰色点云表示原始 ATL03 光子，红色轨迹表示连续性约束后的参考表面。}
\end{figure}

得到参考表面后，需要将不同沿轨位置处的光子对齐到同一垂向坐标系。设 \(h_s(t_i)\) 为第 \(i\) 个光子所在位置的参考表面高程，则表面对齐后的相对高程定义为
\begin{equation}
    h'_i = h_i - h_s(t_i),
\end{equation}
其中 \(h'_i\) 表示第 \(i\) 个光子相对于局地主表面的高程，\(h_i\) 为该光子的原始高程，\(h_s(t_i)\) 为光子所在时刻 \(t_i\) 对应的参考表面高程。经过该变换后，主表面附近对应 \(h'_i=0\)，表面上方主要为大气背景区，表面下方则对应雪层内部多次散射引起的延迟回波。为与后续路径长度统计保持一致，进一步定义深度坐标为
\begin{equation}
    z_i = -h'_i.
\end{equation}
其中 \(z_i\) 为第 \(i\) 个光子相对于主表面的深度坐标。这样，\(z_i>0\) 表示表面以下，\(z_i<0\) 表示表面以上。表面对齐的作用在于消除沿轨表面起伏对垂向剖面的展宽影响，使不同位置处的次表层光子能够在同一深度坐标下进行累积统计。

\subsubsection{时间--深度剖面构建与沿轨聚合}
在统一深度坐标下，本文将光子投影到规则的时间--深度网格中。设第 \(k\) 个时间分箱和第 \(j\) 个深度分箱对应的光子计数为 \(B_{j,k}\)，则 \(B\) 构成沿轨二维光子计数矩阵。随后以长度为 \(\Delta T\) 的滑动窗口沿轨移动，并对窗口内的若干光子列进行平均，得到第 \(m\) 条局地剖面：
\begin{equation}
    S_m(z_j)=\frac{1}{N_m}\sum_{k\in W_m} B_{j,k},
\end{equation}
其中 \(S_m(z_j)\) 为第 \(m\) 条局地剖面在第 \(j\) 个深度分箱 \(z_j\) 处的平均光子计数，\(B_{j,k}\) 为第 \(k\) 个时间分箱和第 \(j\) 个深度分箱中的光子计数，\(W_m\) 为第 \(m\) 个沿轨窗口包含的时间分箱集合，\(N_m\) 为该窗口内有效光子列数量。窗口长度控制剖面的信噪比，前进步长控制沿轨剖面的采样密度；当窗口过短时，次表层弱信号易受随机光子波动影响，当窗口过长时，又会削弱沿轨变化特征。因此，剖面构建需要在垂向分辨率、沿轨分辨率和信号稳定性之间进行折中。

每条剖面同时保存对应窗口内光子的平均经纬度、时间或有效光子数量等辅助信息，用于后续沿轨定位和外部雪深数据配对。剖面统计完成后，表面以上远离主峰的深度区间可用于估计局地背景噪声，表面以下的正深度区间则保留为多次散射信号分析对象。由此得到的 \(S_m(z)\) 不再是原始点云，而是相对于局地主表面对齐、经过沿轨聚合的后向散射剖面。后续系统响应去卷积、路径长度分布恢复以及雪深统计矩反演均基于该剖面展开。

表面对齐完成后，离散光子被投影到 \(-15\sim30\)\,m 的深度范围内，并在沿轨方向通过滑动窗口进行聚合。该处理在物理上相当于用局部样本估计雪层回波强度随深度衰减的包络，使弱长尾信号能够从随机光子波动中显现出来。如\autoref{fig:profile-denoise-dates}所示，噪声抑制后表面以上背景被压低，表面以下仍保留从主峰向深部衰减的长尾结构，说明剖面构建能够兼顾噪声控制与有效信号保留。

\begin{figure}[!htbp]
    \centering
    \includegraphics[width=0.78\textwidth]{final/processing/0419_Day_profile_strg3_ATL03_20190419131031_03250304_007_01.profile.jpg}
    \caption{\label{fig:profile-denoise-dates}2019 年 4 月 19 日典型 ATL03 轨段的时间--深度剖面与噪声抑制结果。红色曲线表示背景扣除前的标准化剖面，蓝色曲线表示背景扣除与平滑后的有效剖面。}
\end{figure}
% \FloatBarrier

\subsection{背景噪声扣除与系统响应校正}
\label{sec:background-response-correction}
ATL03 剖面中的表面以下信号并不全部来自真实雪层散射。太阳背景、随机暗计数、后脉冲和系统响应拖尾都可能污染次表层长尾。如果不进行校正，仪器伪信号可能被误认为雪层内部多次散射，从而造成雪深偏差。因此，本文在路径长度统计之前先进行背景噪声扣除和系统响应去卷积。

背景噪声利用表面以上、远离主峰的区间估计。本文采用 \(-15\sim-8\)\,m 的表面以上区间作为局地背景区，并将该区间内的平均计数作为局地背景水平 \(\hat{b}_m\)。对第 \(m\) 条剖面，背景扣除可写为
\begin{equation}
    S_{b,m}(z)=S_m(z)-\hat{b}_m.
\end{equation}
其中 \(S_{b,m}(z)\) 为第 \(m\) 条剖面扣除背景后的光子计数剖面，\(S_m(z)\) 为扣除前剖面，\(\hat{b}_m\) 为该剖面的局地背景计数估计，\(z\) 为深度坐标。该处理用于削弱太阳背景和随机噪声对次表层长尾信号的污染。背景扣除后，表面以下 \(z>0\) 区域保留为后续反演使用的后向散射剖面；若扣除后剖面整体过弱或有效深度范围内缺少正值，则不再进入后续统计矩计算。
\begin{figure}[!htbp]
    \centering
    \includegraphics[width=0.78\textwidth]{final/method/detector_response_model}
    \caption{\label{fig:detector-response-model}单光子探测器响应模型示意图}
\end{figure}
系统响应校正用于削弱后脉冲和仪器脉冲拖尾对浅层长尾的影响。单光子探测器响应模型如\autoref{fig:detector-response-model}所示。在背景项已扣除后，设观测剖面为 \(S_o(z)\)，校正后剖面为 \(S_c(z)\)，系统响应函数为 \(R(z)\)，则观测过程可写为
\begin{equation}
    S_o(z)=R(z)\otimes S_c(z)+\varepsilon(z),
\end{equation}
其中 \(S_o(z)\) 为观测剖面，\(R(z)\) 为系统响应函数，\(S_c(z)\) 为校正后剖面，\(\otimes\) 表示卷积运算，\(\varepsilon(z)\) 为观测噪声项，\(z\) 为深度坐标。离散形式下，可将卷积过程表示为
\begin{equation}
    \mathbf{F}\mathbf{S}_c=\mathbf{S}_o,
\end{equation}
其中 \(\mathbf{F}\) 为由系统响应函数构成的下三角卷积矩阵，\(\mathbf{S}_o\) 和 \(\mathbf{S}_c\) 分别表示观测剖面与校正后剖面。

去卷积不是普通平滑，而是连接观测剖面与物理反演的必要校正。由于去卷积属于病态线性反问题，在低信噪比条件下可能放大局部噪声，并在部分深度位置产生负值。因此，去卷积结果需要结合非负约束进行筛选；由于后向散射强度在物理上应为非负量，负值不再参与后续统计矩计算。后文将去卷积后用于反演的校正后向散射剖面记为 \(\beta(z)\)。

\begin{figure}[!htbp]
    \centering
    \includegraphics[width=0.78\textwidth]{final/processing/0419_deconv_Day_profile_strg3_ATL03_20190419131031_03250304_007_01.profile.jpg}
    \caption{\label{fig:process-deconv-example}2019 年 4 月 19 日典型轨段的去卷积前后剖面对比。蓝线为观测剖面 \(S_o\)，橙线为校正后剖面 \(S_c\)，紫线为系统响应参考。}
\end{figure}

如\autoref{fig:process-deconv-example}所示，去卷积后浅层拖尾明显收缩，校正后剖面在 \(0\sim5\,\text{m}\) 深度范围内衰减更快，说明系统响应去卷积能够有效削弱表面附近的展宽影响。同时，较深部分的整体衰减包络仍被保留下来，表明当前参数设置在削弱浅层伪信号的同时，没有明显破坏后续路径长度统计所依赖的主要剖面形态。

\FloatBarrier

\subsection{雪深迭代求解}
\label{sec:snow-iteration-parameters}
经过背景扣除和系统响应校正后，每条有效剖面被记为 \(\beta_m(z)\)，其表面以下长尾可用于路径长度统计矩计算。雪深求解流程如\autoref{fig:snow-depth-flowchart}所示，核心步骤包括吸收校正、路径长度矩计算、一阶矩雪深估计、扩散散射参数更新和吸收系数迭代。

\begin{figure}[!htbp]
    \centering
    \includegraphics[
        width=0.9\textwidth,
        height=0.8\textheight,
        keepaspectratio
    ]{final/method/snow_depth_inversion_flowchart}
    \caption{\label{fig:snow-depth-flowchart}基于代码实现的雪深反演流程图}
\end{figure}
% \FloatBarrier

\subsubsection{吸收校正与统计矩计算}
由于较长路径光子在雪层内更容易被吸收，直接对 \(\beta_m(z)\) 积分会低估长尾贡献。本文在有效深度范围内引入 \(e^{2k_az}\) 作为吸收校正因子，并在离散深度网格上计算路径长度统计矩：
\begin{equation}
    \langle L^n\rangle_m =
    \frac{\sum_{z_j\in Z}(2z_j)^n\beta_m(z_j)e^{2k_az_j}\Delta z}
    {\sum_{z_j\in Z}\beta_m(z_j)e^{2k_az_j}\Delta z},
\end{equation}
其中 \(\langle L^n\rangle_m\) 为第 \(m\) 条剖面的路径长度 \(n\) 阶矩，\(z_j\) 为第 \(j\) 个深度分箱中心，\(Z\) 为进入统计的表面以下有效深度集合，\(\beta_m(z_j)\) 为第 \(m\) 条校正后向散射剖面在 \(z_j\) 处的强度，\(k_a\) 为吸收系数，\(\Delta z\) 为垂向分辨率。本文以一阶矩作为雪深估计的主控量，即
\begin{equation}
    H_m=\frac{\langle L\rangle_m}{2}.
\end{equation}
其中 \(H_m\) 为第 \(m\) 条剖面的雪深估计值，\(\langle L\rangle_m\) 为该剖面的路径长度一阶矩。二阶矩不直接作为最终雪深输出，而用于约束扩散散射系数和吸收系数更新。

\subsubsection{吸收系数与扩散散射系数迭代}
为了使路径长度统计与雪层光学参数保持一致，本文在计算过程中采用迭代方式更新吸收系数和扩散相关参数。首先给定吸收系数初值 \(k_a^{(0)}\)，利用 \(\beta_m(z)\) 计算吸收校正后的低阶矩，并由一阶矩得到初始雪深。随后根据二阶矩估计扩散散射系数：
\begin{equation}
    k_{sd}=\frac{\langle L^2\rangle_m}{H_m^3}.
\end{equation}
在获得反照率 \(a\) 后，可进一步估计等效雪粒径
\begin{equation}
    R=\frac{1}{k_a}\left(\frac{1-a}{8.43}\right)^2,
\end{equation}
并计算扩散衰减系数
\begin{equation}
    k_d=0.65\sqrt{\frac{k_a}{R}}.
\end{equation}
据此更新吸收系数
\begin{equation}
    k_a'=\frac{k_d^2}{3(k_{sd}+k_a)}.
\end{equation}
其中 \(k_{sd}\) 表示扩散散射系数，\(\langle L^2\rangle_m\) 为第 \(m\) 条剖面的路径长度二阶矩，\(R\) 为等效雪粒径，\(a\) 为有效反照率，\(k_d\) 为扩散衰减系数，\(k_a'\) 为更新后的吸收系数，式中的 8.43、0.65 和 3 为该参数化关系中的经验或理论常数。当相邻两次迭代的吸收系数差小于设定阈值时，认为雪层吸收、散射参数与路径长度矩之间达到一致，并输出对应雪深。

为了展示上述统计矩求解在沿轨方向上的输出形式，本文进一步给出 2019 年 4 月 19 日典型轨段的雪深反演结果。图中灰线表示逐廓线雪深，红线表示 9 点滑动平均结果，用于削弱单个窗口内光子计数波动造成的局部尖峰。由\autoref{fig:track-snow-depth-three-days}可见，平滑曲线基本保留了逐廓线雪深的沿轨变化趋势，同时抑制了少量突变点，说明滑动窗口聚合、背景扣除、系统响应校正和统计矩反演能够形成连续的沿轨雪深序列。

\begin{figure}[H]
    \centering
    \includegraphics[width=0.92\textwidth]{final/processing/track_plots/0419_ResultDay_profile_strg3_ATL03_20190419131031_03250304_007_01.profile_track.png}
    \caption{2019 年 4 月 19 日典型 ATL03 轨段的沿轨雪深反演结果。灰线表示逐廓线雪深，红线表示 9 点滑动平均雪深。}
    \label{fig:track-snow-depth-three-days}
\end{figure}
\FloatBarrier

\subsection{本章小结}
本章建立了从 ATL03 光子事件到雪深结果的完整方法链。首先根据主表面识别和表面对齐构建统一深度坐标，再通过背景扣除和系统响应去卷积获得校正后的后向散射剖面，最后利用多次散射路径长度统计矩和吸收系数迭代求解沿轨雪深。该方法的关键环节包括表面识别、剖面构建、背景扣除、去卷积和统计矩求解。第4章将分别利用 OIB 海冰验证数据以及 UA/CMC 陆地验证数据对反演结果进行评价。
